\chapter{Frailty}

\section*{Definition}
Frailty is a geriatric syndrome of decreased physiologic reserve and increased vulnerability to stressors, clinically defined by the presence of \ensuremath{\geq}3 of 5 Fried criteria: unintentional weight loss, exhaustion, weak grip strength, slow gait speed, and low physical activity. It predicts adverse outcomes including falls, hospitalization, disability, and mortality.

\section*{What Happens in Your Body}
Frailty arises from cumulative age-related decline across multiple organ systems, particularly musculoskeletal, neuroendocrine, and immune axes. Dysregulation of inflammatory immune signaling proteins (IL-6, TNF-\ensuremath{\alpha}), sarcopenia (loss of muscle mass and strength), and mitochondrial dysfunction create a vicious cycle of declining energy metabolism and functional capacity. Reduced growth hormone and testosterone levels, along with insulin resistance, further accelerate catabolism.

\section*{Causes}
\begin{itemize}
  \item \textbf{Aging-related:} Sarcopenia, osteopenia, immunosenescence, neurodegenerative changes
  \item \textbf{Chronic disease:} Heart failure, COPD, diabetes mellitus, chronic kidney disease, osteoarthritis
  \item \textbf{Nutritional:} Protein--energy malnutrition, vitamin D deficiency, micronutrient deficits
  \item \textbf{Psychosocial:} Depression, social isolation, cognitive impairment, polypharmacy
  \item \textbf{Acute triggers:} Hospitalization, surgery, infection, injury
\end{itemize}

\section*{When to See a Doctor}
See your doctor if:
\begin{itemize}
  \item Unintentional weight loss of \ensuremath{>}5\% in 6--12 months or decline in functional status
  \item Recurrent falls, new mobility limitation, or inability to perform activities of daily living
  \item Post-hospitalization deconditioning or failure to return to baseline
\end{itemize}

\section*{Related Symptoms}
\begin{itemize}
  \item Sarcopenia, osteopenia/sarcopenic obesity, slow gait, weakness, fatigue, weight loss, poor balance
\end{itemize}
\textbf{Important:} Frailty is potentially reversible with early, targeted interventions, making routine screening in adults \ensuremath{\geq}70 years a recommended component of geriatric care.

\section*{Self-Care / Home Management}
\begin{itemize}
  \item Engage in progressive resistance training and balance exercises (e.g., Tai Chi, physical therapy) at least 2--3 times per week
  \item Consume adequate protein (1.2--1.5 g/kg/day) and vitamin D (800--1000 IU/day)
  \item Review medications with a pharmacist to deprescribe unnecessary or high-risk drugs
\end{itemize}

\section*{How Your Doctor Will Evaluate You}
\begin{itemize}
  \item Screen with FRAIL scale, Fried phenotype, or Clinical Frailty Scale
  \item Assess gait speed (<0.8 m/s is abnormal), grip strength (dynamometry), and 5-times sit-to-stand time
  \item Evaluate for reversible contributors: depression screening, thyroid function, vitamin B12, vitamin D, anemia workup
\end{itemize}

\section*{Treatment Options}
\begin{itemize}
  \item Multicomponent exercise program combining resistance, aerobic, and balance training
  \item Nutritional optimization: protein supplementation, oral nutritional supplements if underweight
  \item Medication reconciliation and deprescribing of anticholinergics, sedatives, and antihypertensives
  \item Comprehensive geriatric assessment and multidisciplinary case management for frequent fallers or recurrent hospitalizations
\end{itemize}

\section*{References}
\begin{thebibliography}{99}
\bibitem{frailty:ref1} Fried LP, Tangen CM, Walston J, et al. ``Frailty in Older Adults: Evidence for a Phenotype.'' J Gerontol A Biol Sci Med Sci. 2001;56(3):M146--M156.
\bibitem{frailty:ref2} Clegg A, Young J, Iliffe S, et al. ``Frailty in Elderly People.'' Lancet. 2013;381(9868):752--762.
\bibitem{frailty:ref3} Morley JE, Vellas B, van Kan GA, et al. ``Frailty Consensus: A Call to Action.'' J Am Med Dir Assoc. 2013;14(6):392--397.
\bibitem{frailty:ref4} Walston J, Hadley EC, Ferrucci L, et al. ``Research Agenda for Frailty in Older Adults: Toward a Better Understanding of Physiology and Etiology.'' J Am Geriatr Soc. 2006;54(6):991--1001.
\bibitem{frailty:ref5} Dent E, Morley JE, Cruz-Jentoft AJ, et al. ``Physical Frailty: ICFSR International Clinical Practice Guidelines for Identification and Management.'' J Nutr Health Aging. 2019;23(9):771--787.
\bibitem{frailty:ref6} Xue QL. ``The Frailty Syndrome: Definition and Natural History.'' Clin Geriatr Med. 2011;27(1):1--15.
\end{thebibliography}
